\section{Frobenius traces and compact Weil classes}\label{tr:section}

This section proves the concentration and spanning assertions needed
for the image definition of $\cF_Y$.  They will also permit every
function in the second variable of the automorphic pairing to be
represented by a finite combination of compact spectral tests.
The argument has three parts: remove invariant-base support on the
trace, contract the nonconstant part of an orbit trace, and use finite
supported localization to assemble the result.

\subsection{The invariant-base projector and actual supported lifts}

\begin{lemma}[Invariant support on traces]\label{tr:invariant-support}
For every Frobenius-stable conical $z$-support bound $\Omega$, inclusion
induces an isomorphism
\begin{equation}\label{tr:remove-completion}
 \Tr(\Phi,\operatorname{Ind}\mathscr B_{\Omega,0})
       \xrightarrow{\;\sim\;}
 \Tr(\Phi,\operatorname{Ind}\mathscr B_\Omega).
\end{equation}
A compact Weil object on the right has a compact Weil lift on the left
whose trace class maps to a nonzero scalar multiple of its class.
Both statements hold for the compact opposite categories.
Moreover, for Hom pairings of two such lifts, the finite Koszul
construction multiplies each orbit contribution by a constant nonzero
factor, independent of the compact stabilizer variable.
\end{lemma}

\begin{proof}
The ideal $I\subset R^J$ of the invariant origin is generated by a
finite Frobenius-stable homogeneous space $U$ of positive-degree
invariant polynomials.  To obtain $U$, choose homogeneous algebra
generators and take their Frobenius spans inside the finitely many
finite-dimensional polynomial-degree spaces involved.  The inverse
Hom action has strictly negative weights on $U$ by
Lemma~\ref{mod:hom-action}; its forward inverse has strictly positive
weights.  In either convention the appropriate transport operator
$F_U$ has no eigenvalue one.  Thus $1-F_U$ is invertible, even when
$F_U$ has Jordan blocks.

For a central scalar $a$, its two actions on a twisted Hochschild
complex, one transported by Frobenius, are chain homotopic.  This is
the elementary homotopy moving $a$ once around the bar.  Applying it
simultaneously to a basis of $U$ and multiplying by $(1-F_U)^{-1}$
shows that every generator of $I$ acts trivially on trace cohomology.
Equivalently, regard the complex as a module over the polynomial
algebra $\Sym U$ mapping to $R^J$; its cohomology is supported at the
origin.  In particular it is $I$-torsion, a statement about the ordinary
complex rather than about its completion.

Let $i$ be inclusion of the support category and let $\Gamma_I$ be
its continuous colocalization.  On the ambient category,
$i\Gamma_I$ is tensor product with the scalar local-cohomology complex
$\RGamma_I(R)$.  The inclusion preserves compact objects.  Cyclicity
of categorical trace, followed by continuity in its coefficient
bimodule, therefore identifies the map in
\eqref{tr:remove-completion} with
\[
 \RGamma_I\Tr(\Phi,\operatorname{Ind}\mathscr B_\Omega)
       \longrightarrow
 \Tr(\Phi,\operatorname{Ind}\mathscr B_\Omega).
\]
One can compute this equality directly with the finite Koszul
complexes on powers of the chosen equations and their filtered union.
It is an isomorphism because the target cohomology is $I$-torsion.
This proves the trace assertion without replacing Hochschild chains
by a completed bar complex.

To lift an actual class, let $M$ be a compact Weil object and let
$K(U)$ be the finite equivariant Koszul complex on the map
$U\otimes R\to R$.  Then $M\otimes K(U)$ is compact, has ordinary
support over $I=0$, retains its $z$-support, and has the induced Weil
structure.  The finite exterior filtration and additivity of trace
give
\begin{equation}\label{tr:koszul-factor}
 [M\otimes K(U)]
     =\left(\sum_j(-1)^j\Tr(F_U\mid\Lambda^jU)\right)[M]
     =\det(1-F_U\mid U)\,[M].
\end{equation}
The determinant is nonzero.  The notation $F_U$ in this formula means
the coefficient transport on the chosen exterior terms; passing to
inverse Weil structures replaces it by its inverse, which still has
no eigenvalue one.  Dividing the class by this nonzero complex scalar
gives any desired lift in the trace vector space, with the representative
itself still an actual compact Weil object.

The same finite filtration proves the pairing statement.  In a Hom
from one Koszul lift to another, its associated terms have coefficient
factor $\Hom(\Lambda^aU,\Lambda^bU)$, with the dual transport on the
contravariant coefficient and the original transport on the covariant
one.  Its alternating character is the product of the corresponding
two nonzero determinants.  Since the equations are $J$-invariant,
these are constant functions of every compact $L_e$-variable.
Taking sections with support in an orbit commutes with this finite
filtration.  Thus the same factor multiplies each orbit contribution,
not merely the sum of all contributions.  This also proves the
assertions for opposite categories and their inverse structures.
\end{proof}

\subsection{An algebraic contraction of the twisted bar complex}

We first identify the original $\Phi$-graph with the transported model,
then use the contracting inverse action $S$ on that same graph.
On a finite stable list of models let $\pi$
be the permutation induced by Frobenius and choose
$\alpha_i:\Phi M_i\xrightarrow{\sim}M_{\pi i}$.
These identifications give a functor $\Psi$ on the chosen models.
For arrows $a:M_i\to M_j$ and $b:M_{\pi i}\to M_{\pi j}$,
its forward transport and inverse are
\begin{equation}\label{tr:forward-inverse}
 \begin{split}
 T(a)=\Psi(a)&=\alpha_j\Phi(a)\alpha_i^{-1},\\
 S(b)=T^{-1}(b)&=\Phi^{-1}(\alpha_j^{-1}b\alpha_i).
 \end{split}
\end{equation}
The action $S$ is the contracting Hom action of
Lemma~\ref{mod:hom-action}.

To match trace classes $M\to\Phi M$, use the left-twisted graph
$G_\Phi(i,j)=\Hom(M_i,\Phi M_j)$, whose left action is
composition through $\Phi$.  The maps
\begin{equation}\label{tr:graph-identification}
 \theta_{ij}:G_\Phi(i,j)\longrightarrow G_\Psi(i,j),
 \qquad m\longmapsto\alpha_j m,
 \qquad G_\Psi(i,j)=\Hom(M_i,M_{\pi j})
\end{equation}
identify the graph bimodules: for composable arrows,
$\theta(\Phi(a)mb)=\Psi(a)\theta(m)b$, with the appropriate
source and target indices.  Transport the entire twisted bar
construction through this identification.  In the $G_\Psi$
coordinates simultaneous forward transport is $\Psi$ on every
factor.  On the original graph coefficient it is instead
\begin{equation}\label{tr:graph-transport}
 T_G(m)=\Phi(\alpha_j)\Phi(m)\alpha_i^{-1}
          \in G_\Phi(\pi i,\pi j),
 \qquad
 \theta_{\pi i,\pi j}(T_G(m))=\Psi(\theta_{ij}(m)).
\end{equation}
For example, for one algebra with automorphism $\phi$ and
$\psi=\operatorname{Ad}_\alpha\phi$, these formulas are
$\theta(m)=\alpha m$ and
$T_G(m)=\phi(\alpha)\phi(m)\alpha^{-1}$.
Thus the original $\Phi$-trace is retained, with its graph coefficient
canonically identified.  Ordinary-arrow transport cannot be applied
to the raw graph factor without this identification.

\begin{lemma}[Locally finite contraction]\label{tr:bar-contraction}
Let a small generating differential graded category with an
autoequivalence be modeled in complexes of locally finite cyclic
representations.  Suppose, after finite cyclic permutations of its
labels and pure choices of their structures, its morphism cohomology
has nonpositive absolute $S$-weights; composition adds weights.
Suppose its weight-zero morphism cohomology is concentrated in degree
zero and its additive idempotent completion is semisimple, with
endomorphism algebra $\C$ for every simple object.  Then its twisted
trace is the twisted trace of that split semisimple category.  It is concentrated in
degree zero, with one basis line for each fixed simple label.
Nontrivially permuted simple labels contribute zero.
\end{lemma}

\begin{proof}
Take the sum of the nonpositive-weight subcomplexes in every morphism
complex.  This is a differential graded subcategory: identities have
weight zero, the differential preserves weights, and composition adds
them.  It is quasi-equivalent to the original category because the
excluded weight pieces have zero cohomology.  Locally finite
semisimple-weight decomposition is exact, so this assertion also holds
for unbounded morphism complexes.

Use the direct-sum realization of its twisted cyclic bar complex.
Every chain belongs to a finite sum of finite-dimensional cyclic
subrepresentations.  On a finite permutation-stable list of labels,
take a common power fixing the labels; divided absolute weights for
that power define the same nonpositive grading on the bar terms.
Only the weights are read using this power: the functor and graph
whose trace is computed remain $\Phi$ and $G_\Phi$.
In the identified left-graph convention, cyclic rotation has the
one-algebra formula
$t_n(a_0\mid\cdots\mid a_n)
 =(\Psi(a_n)\mid a_0\mid\cdots\mid a_{n-1})$,
with the usual Koszul signs for graded factors.
Its $(n+1)$-fold iterate applies $\Psi$ to every factor.
The usual paracyclic homotopy therefore gives
\begin{equation}\label{tr:paracyclic}
 dH+Hd=1-T.
\end{equation}
Here $d$ is the total Hochschild and internal differential, and the
operator $T$ is simultaneous forward transport in the identified
$G_\Psi$ coordinates, or equivalently the ordinary-arrow transport
and \eqref{tr:graph-transport} on the original graph.  The identity can be obtained
by inserting the unit and summing the cyclic rotations: consecutive
faces cancel in pairs, and the two remaining faces are the identity
and a full rotation.  Transporting this identity through
\eqref{tr:graph-identification} gives the asserted identity on the
original $\Phi$-twisted bar.  In degree zero the one-algebra convention
makes it explicit: with $d(m\mid a)=ma-\phi(a)m$, the term
$H_0(m)=\alpha^{-1}\mid\alpha m$ satisfies
$dH_0(m)=m-T_G(m)$.

Since $T$ commutes with $d$, the operator
\begin{equation}\label{tr:inverse-homotopy}
 H_{\mathrm{inv}}=-T^{-1}H
 \qquad\text{satisfies}\qquad
 dH_{\mathrm{inv}}+H_{\mathrm{inv}}d=1-S.
\end{equation}
On every strictly negative-weight block, $1-S$ is invertible.
This inverse is algebraic: on each finite-dimensional cyclic
subrepresentation it is a polynomial in $S$, by its minimal
polynomial, and the resulting inverses agree on inclusions.
There is no infinite geometric series of chains.
The homotopy $(1-S)^{-1}H_{\mathrm{inv}}$ contracts that block.
The homotopy preserves total weight, and the same conclusion holds
when a power is used to read the weights, since every eigenvalue of
$S$ then has modulus strictly less than one.

Only weight zero remains.  As every arrow has nonpositive weight, a
weight-zero bar term has only weight-zero arrows.  Hence this is
exactly the bar complex of the weight-zero category.  Its morphism
cohomology is in degree zero; the canonical differential graded
truncation gives a zigzag of quasi-equivalences to its degree-zero
cohomology category, compatibly with the cyclic action and composition.
A split semisimple category has the same trace as its full subcategory
of simple objects.  By the scalar-endomorphism hypothesis, its only
nonzero morphisms are scalar endomorphisms.  The twisted bar therefore gives $\C$ in degree
zero for a fixed simple and zero for a permuted block.

Finally, any chain uses only finitely many labels.  Enlarge them to a
finite permutation-stable list and apply the preceding argument.
Filtered union gives the assertion for all labels.  All realizations
are direct sums; neither an infinite product nor convergence of an
infinite Hochschild spectral sequence has been used.
\end{proof}

\begin{proposition}[Trace of an orbit]\label{tr:orbit-trace}
The trace on the supported category of a Frobenius-stable nilpotent
orbit is concentrated in degree zero.  A basis is given by the classes
of $M_P$ for Frobenius-fixed simple $P\in\mathsf E_e$, with chosen
pure normalized structures.  Every basis class lifts to an actual
compact Weil object supported on the orbit closure, and also to one
with invariant null-fiber support up to a nonzero scalar.
The same conclusions hold for the compact opposite category with
simultaneous Frobenius.
\end{proposition}

\begin{proof}
By Proposition~\ref{orb:generation}, the full subcategory on the simple
labels generates the supported category; Morita invariance permits
computing its trace on this generating subcategory.  We first check
the local finiteness required by Lemma~\ref{tr:bar-contraction}, since
numerical weight bounds on cohomology alone would not suffice.

All models used in Section~\ref{orb:section} are algebraic equivariant
models with finite-dimensional starting representation data.  Keep
the grading internal until direct-sum totalization.  For a finite
cyclic list of labels, take a common finite-type algebraic cover $H$
of the equivariance and cyclic chain-map data, acting on the chosen
objects and the invariant quasi-compact
separated ambient open $U$.  Polynomial functions, finite Koszul
resolutions, equivariant restriction, and proper pushforward are
defined in rational representations of this group.

One can compute derived sections without choosing a
transporter-stable affine cover.  Equivariant derived pushforward
along $[U/H]\to BH$ produces complexes of rational
$H$-representations.  Flat base change along $\operatorname{pt}\to BH$
recovers ordinary derived sections on $U$, and the lax monoidal
pushforward supplies their composition maps
\cite[\S1.3.10, Corollary~1.4.5, and \S\S6.4.3--6.4.6]{DGFiniteness}.
Here differential Hom is an ordinary square-zero complex because
the two curvature terms cancel.  Functorial bar--cobar rectification
uses direct sums of finite tensor words, so it stays in rational
representations, preserves the cyclic action, and gives a compatible
differential graded model with strict composition
\cite[Proposition~2.1 and \S\S4.3, 4.8, 5.2]{KellerAInfinity}.
Restriction to the algebraic closure of the cyclic operator makes
the action on this model locally finite.  Stabilizer invariants are
computed by the finite Lie-cochain complex followed by exact
reductive invariants.  These bounded constructions commute with the
direct-sum folding of the grading.  Thus the locally finite cyclic
actions are defined on morphism complexes with their composition,
and decomposition by the characters of the semisimple part gives
the weight subcomplexes used in the contraction lemma.

Proposition~\ref{orb:weil-lifts} provides finite cyclic orbits of the
labels and consistently pure implementing maps, including when a
power is required to fix a label.  On taking invariants, the
triple translation is by a group element and does not change the
actual arrow action $S$.  Proposition~\ref{orb:hom-calculation}
therefore says that this model has no positive-weight cohomology and
that its zero-weight category is exactly the semisimple category
$\mathsf E_e$.  Its simple objects have scalar endomorphisms: each is
a finite-dimensional complex module, so Schur's Lemma applies over
the algebraically closed field $\C$.
Lemma~\ref{tr:bar-contraction} proves concentration
and the stated basis.  The incidence construction of
Proposition~\ref{orb:weil-lifts} supplies its actual compact
orbit-closure representatives; Lemma~\ref{tr:invariant-support}
supplies the completed representatives.

For the opposite category use the same object models, with reversed
arrows and inverse object structures.  Its simultaneous forward
action on a reversed arrow is the forward action on the underlying
original arrow; its inverse is consequently the same contracting
$S$ on that underlying complex.  The weight calculation and
paracyclic identity are unchanged.  In particular, this argument
uses the autoequivalence $\Phi^{\mathrm{op}}$ of the opposite
category, not the inverse autoequivalence of the original category.
It proves the same spanning assertion with actual inverse Weil
structures.
\end{proof}

\subsection{Finite support filtrations and the full function space}

\begin{theorem}[Concentration, injection, and Weil spanning]\label{tr:spanning}
On a fixed retained spectral component, the trace of every closed
invariant nilpotent support category is concentrated in degree zero.
If $\Omega'\subset\Omega$ are such support bounds, then
\begin{equation}\label{tr:support-injection}
 H^0\Tr(\Phi,\operatorname{Ind}\mathscr B_{\Omega',0})
   \longrightarrow
 H^0\Tr(\Phi,\operatorname{Ind}\mathscr B_{\Omega,0})
\end{equation}
is injective.  A finite filtration by invariant unions of orbits has
the orbit traces as its actual successive quotients.  Every trace
class is a finite linear combination of classes of compact objects
with Weil structures; those objects can be chosen with support in the
prescribed bound.  All assertions hold on the compact opposite side.

After assembling retained components, these statements apply to
$\cC_Y$ and $\cC$.  In particular, the image defining $\cF_Y\cA_{\E}$
is the injective image of the supported trace, and every function in
$\cA_{\E}$ and every vector in $\cF_Y\cA_{\E}$ is a finite linear
combination of the corresponding actual compact Weil trace functions.
The analogous statement supplies all tests through the
nonstandard self-duality of Theorem~\ref{inp:duality}.
\end{theorem}

\begin{proof}
Choose a finite filtration of the nilpotent support by closed invariant
unions of $J$-orbits, grouping Frobenius permutations into stable strata.
It exists because there are finitely many nilpotent orbits and
Frobenius preserves their closure order.  At each step,
Lemma~\ref{orb:localization} gives a compactly generated localization.
Twisted Hochschild homology is localizing, so it gives a cofiber
sequence of traces; this is the trace-localization formalism of
\cite[Theorem~3.4 and Proposition~5.4]{HSS}, applied to the actual
Frobenius-compatible localization just constructed.
Disjoint orbits closed in the corresponding open
have orthogonal supported categories.  If Frobenius permutes such
summands without a fixed summand, their twisted trace is zero: in the
bar model, nonzero chains would have to begin and end in the same
summand, while the graph coefficient ends in its Frobenius translate.
Every fixed orbit has the degree-zero trace of
Proposition~\ref{tr:orbit-trace}.

Induction up this finite filtration proves concentration in degree
zero.  Each cofiber sequence is therefore a short exact sequence on
$H^0$.  More generally, apply the same argument to the finite orbit
filtration of the complement of $\Omega'$ in $\Omega$.  Its trace is
also concentrated in degree zero, so the localization sequence starts
\[
 0\longrightarrow H^0\Tr(\Phi,\mathscr B_{\Omega',0})
 \longrightarrow H^0\Tr(\Phi,\mathscr B_{\Omega,0})
 \longrightarrow H^0\Tr(\Phi,\mathscr B_{\Omega,0}/
                                      \mathscr B_{\Omega',0})
 \longrightarrow0,
\]
where taking ind-completions in the trace notation is understood.
This proves \eqref{tr:support-injection} and identifies the successive
quotients with the computed orbit traces.  It does not infer
injectivity from concentration of the full category alone.

At an induction step, each quotient basis class lifts by the compact
incidence construction and the finite invariant Koszul construction.
Subtract a finite linear combination of these lifts from a given
class; the remainder lies in the preceding supported trace.
The induction terminates after finitely many strata.  Thus every
class is spanned by actual compact Weil classes with the required
closed support.  This argument also proves spanning on the opposite
side; its localization sequences and orbit lifts have already been
checked in Proposition~\ref{tr:orbit-trace}.

For the retained union of components, compact objects involve only
finitely many components and there are no morphisms between distinct
components.  The same bar argument shows
\begin{equation}\label{tr:component-sum}
 \Tr(\Phi,\cC)
   =\bigoplus_{\substack{Z_a\subset Z\text{ retained}\\
                         \Phi Z_a=Z_a}}
          \Tr(\Phi,\cC|_{Z_a}),
\end{equation}
and likewise for supported categories and their compact opposites.
A component permuted nontrivially contributes zero even if a power
fixes it.  The direct sum in \eqref{tr:component-sum} is essential:
no vector acquires infinitely many component contributions.

The Frobenius-compatible restricted equivalence of
Theorem~\ref{inp:langlands} transports these results to the categories
in the problem, retaining exactly the prescribed prime union.
The trace/functions isomorphism and ordinary compact local-term
compatibility of Theorem~\ref{inp:trace} take each compact Weil class
to its actual finitely supported trace function.  This identifies the
supported trace injection with the image used to define
$\cF_Y\cA_{\E}$.  Finally the simultaneous-Frobenius nonstandard duality
identifies the trace on the compact opposite category with the full
space of tests.  The opposite spanning statement therefore supplies
every such test by a finite linear combination of the constructed
classes.  All calculations were performed after an arbitrary
coefficient isomorphism with $\C$; concentration, injections, and finite
spanning transport back along that isomorphism.
\end{proof}
